```latex
% ch1_introduction.tex

\chapter{Introduction}
\label{ch:introduction}

Le clustering de grands volumes de données présente des contraintes particulières lorsque les observations arrivent sous forme de flux. Contrairement à une approche batch classique, il n'est pas toujours possible de conserver l'ensemble des données en mémoire ni de les parcourir plusieurs fois. Les algorithmes de clustering adaptés aux flux doivent ainsi limiter leur consommation mémoire tout en maintenant une qualité de partitionnement satisfaisante.

StreamKM++, proposé par Ackermann et al. (2012) \cite{ackermann2012}, répond à cette problématique en construisant progressivement un \emph{coreset}, c'est-à-dire un résumé pondéré et de taille limitée des données observées. Le clustering final est ensuite réalisé sur ce résumé à l'aide de k-means++ \cite{arthur2007}. Cette approche permet de réduire la quantité de données manipulée lors de l'étape finale de clustering tout en conservant une approximation représentative du jeu de données initial.

Bitsakis (2018) \cite{bitsakis2018} a proposé une implémentation distribuée de StreamKM++ avec Apache Flink. Le présent projet étudie l'adaptation de cette approche à Apache Spark, notamment à travers les RDD et Structured Streaming. Il ne s'agit pas uniquement de transposer l'implémentation existante, mais également de l'adapter aux mécanismes propres à Spark, tels que le partitionnement des données, les opérations de réduction distribuée et le traitement par micro-batches.

Le travail s'articule autour de trois objectifs principaux. Le premier consiste à implémenter et à valider les composants algorithmiques nécessaires à StreamKM++, notamment le k-means pondéré, l'initialisation k-means++, la construction du coreset et la stratégie \emph{merge-and-reduce}. Le deuxième objectif concerne leur adaptation à une exécution distribuée avec Spark, où chaque partition construit un résumé local avant la fusion des résultats. Enfin, l'implémentation est évaluée expérimentalement afin d'étudier la qualité du clustering obtenu, l'influence de la taille du coreset ainsi que le comportement de la solution lorsque le volume de données et le niveau de parallélisme augmentent.

L'architecture retenue distingue la logique algorithmique de l'orchestration distribuée. Les packages \texttt{streamkm.core} et \texttt{streamkm.coreset} regroupent les composants indépendants de Spark, tandis que \texttt{streamkm.spark} prend en charge leur exécution distribuée. Cette séparation permet de valider les composants algorithmiques indépendamment de leur distribution avant d'évaluer leur comportement dans un environnement Spark.
```
